% FILE: 00_project_identity.tex
% PROJECT: phd-thesis
% MANUFACTURED: 2026-06-01
% AUTHORITY: phd-thesis/ledgers/WORKFLOW_RECEIPT.yaml (phase: SETUP_COMPLETE)
% DO NOT CLAIM: See ledgers/DO_NOT_CLAIM_LEDGER.md

\section{Project Identity}
\label{sec:phd-thesis-identity}

\subsection{Purpose}
\label{sec:phd-thesis-identity-purpose}

The \texttt{phd-thesis} workspace is the dissertation corpus manufacturing workspace for the
\texttt{process-intelligence} research foundry. Its declared purpose, as stated in
\texttt{phd-thesis/README.md}, is to manufacture the full artifact corpus for a PhD
dissertation grounded in the Van der Aalst Constitution and the full-lifecycle process
intelligence doctrine. Crucially, the README draws an explicit distinction: this workspace does
not produce documents in the ordinary sense --- it manufactures receipted, ledger-backed,
proof-gated artifacts through a declared manufacturing pipeline.

The workspace is the command-and-control index for the dissertation. It is not itself a chapter
source or a PDF; it is the meta-ledger layer from which chapter manufacture is authorized.
Every downstream artifact produced from this workspace is governed by the claim prohibitions
recorded in \texttt{ledgers/DO\_NOT\_CLAIM\_LEDGER.md} and the toolchain provenance recorded
in \texttt{ledgers/WORKFLOW\_RECEIPT.yaml}. No artifact may be released from this workspace
until the ALIVE criteria defined in the README are satisfied: TeX source files exist, the source
ledger is populated, the claim ledger has been respected, the PDF compiles without error, and the
compiled PDF hash is recorded in the receipt chain.

\subsection{Repository Surface}
\label{sec:phd-thesis-identity-surface}

The \texttt{phd-thesis} workspace sits at
\texttt{/Users/sac/process-intelligence/phd-thesis} and operates on branch
\texttt{phd-thesis-corpus-manufacture-001}
(confirmed in \texttt{ledgers/WORKFLOW\_RECEIPT.yaml}). Its directory layout declares six
regions of responsibility:

\begin{itemize}
  \item \texttt{frontmatter/} --- Title page, abstract, acknowledgments, table of contents.
  \item \texttt{chapters/} --- Dissertation chapter \TeX{} sources.
  \item \texttt{projects/} --- Chapter-scoped project artifacts and source evidence.
  \item \texttt{ledgers/} --- Receipt chain, claim ledger, workflow receipt.
  \item \texttt{scripts/} --- Manufacturing pipeline scripts.
  \item \texttt{build/} --- Compiled PDF output and build artifacts.
\end{itemize}

At the time of this writing the \texttt{chapters/}, \texttt{frontmatter/}, \texttt{build/},
and \texttt{scripts/} directories are empty stubs. The populated surface is
\texttt{ledgers/} and \texttt{projects/}. The \texttt{ledgers/PROJECT\_INDEX.yaml}
catalogues 34 research projects across the foundry with their absolute paths, detected
languages, research surfaces, thesis roles, and key files. This makes the workspace the sole
canonical map of the dissertation's evidentiary corpus.

The detected languages recorded in the project manifest are YAML; the detected framework is
\texttt{ggen} (the Governance Generation Engine). The research surfaces are
\texttt{receipt-chain}, \texttt{manufacturing-pipeline}, and \texttt{proof-gate} ---
reflecting the workspace's role as a manufacturing authority rather than a content surface.

\subsection{Primary Research Role}
\label{sec:phd-thesis-identity-role}

The thesis role assigned to this project in \texttt{ledgers/PROJECT\_INDEX.yaml} is
\texttt{dissertation-corpus-manufacturing-workspace}. This distinguishes it from every other
project in the index: while other projects carry roles such as
\texttt{immutable-process-law-foundation} (doctrine/) or
\texttt{core-wasm-execution-authority} (sources/wasm4pm), the \texttt{phd-thesis} project is
the manufacturing authority that orchestrates all of them. It does not contribute research
findings directly; it is the receipted envelope that ensures the research findings are
assembled, cited, and released under the ALIVE criteria.

This role has a formal consequence: the \texttt{phd-thesis} workspace is a gated artifact.
No downstream use of artifacts from this workspace is authorized until the research program
issues authorization --- a constraint stated verbatim in the README and enforced by the
DO\_NOT\_CLAIM\_LEDGER. As of 2026-06-01, the upstream research program has achieved ALIVE
status (PI\_RESEARCH\_PROGRAM\_ALIVE\_001, 12/12 audit gates PASS), which constitutes the
prerequisite authorization for chapter manufacturing to begin.

\subsection{Key Artifacts}
\label{sec:phd-thesis-identity-artifacts}

The key artifacts confirmed present at the \texttt{phd-thesis} workspace are:

\begin{itemize}
  \item \texttt{ledgers/WORKFLOW\_RECEIPT.yaml} --- Manufacturing workflow provenance record.
    Records toolchain availability: pdfTeX 3.141592653-2.6-1.40.26 (TeX Live 2024),
    Python 3.14.3, git 2.51.2, ripgrep 15.1.0. Records phase \texttt{SETUP\_COMPLETE}
    and the pre-manufacture git commit hash
    \texttt{b4279892881e5a8ff83f31b273219a0740fead0b}.
  \item \texttt{ledgers/DO\_NOT\_CLAIM\_LEDGER.md} --- Absolute claim prohibitions:
    no production deployment claims without receipt; no fabricated claims; manufactured
    not ``generated''; no celebrity thought-leader name-dropping; no ``unhackable'';
    ALIVE verdicts require TeX files, source ledger, claim ledger, compiled PDF, and
    recorded PDF hash.
  \item \texttt{ledgers/PROJECT\_INDEX.yaml} --- Canonical ledger of 34 research projects
    with absolute paths, thesis roles, and key files. Manufactured 2026-06-01.
  \item \texttt{projects/phd-thesis/project\_manifest.yaml} --- Manifest declaring the
    workspace's own thesis role, detected languages, frameworks, and key files.
  \item \texttt{README.md} --- Declares the workspace purpose, directory structure,
    ALIVE criteria, and doctrine constraints.
\end{itemize}

The artifacts that are \emph{not yet present} include TeX chapter sources, a compiled PDF,
and a PDF hash receipt. These absences are the reason the workspace holds status PARTIAL
rather than ALIVE. Their manufacture is the declared next phase of this corpus.
